\documentclass[letterpaper,10pt]{article}
\usepackage[empty]{fullpage}\usepackage{titlesec}\usepackage{enumitem}\usepackage[hidelinks]{hyperref}\input{glyphtounicode}\pdfgentounicode=1\pagestyle{empty}
\addtolength{\oddsidemargin}{-0.55in}\addtolength{\textwidth}{1.10in}\addtolength{\topmargin}{-0.65in}\addtolength{\textheight}{1.25in}\raggedright\setlength{\parindent}{0pt}
\titleformat{\section}{\vspace{-4pt}\scshape\large}{}{0em}{}[\titlerule\vspace{-5pt}]
\newcommand{\e}[4]{\textbf{#1}\hfill #2\\\textit{#3}\hfill\textit{#4}\vspace{-3pt}}\newcommand{\bi}{\begin{itemize}[leftmargin=.17in,itemsep=0pt,topsep=2pt]}\newcommand{\ei}{\end{itemize}\vspace{-4pt}}
\begin{document}{\Large\textbf{Giacomo Cappelletto}}\\[-1pt]Boston, MA | giacomo.cappelletto@icloud.com | +1 (857) 753-0133 | \href{https://github.com/JJCAPPE}{GitHub} | \href{https://cv-nu-sage.vercel.app/}{Portfolio}
\section{Education}\e{Boston University}{Boston, MA}{B.S. Computer Engineering, GPA: 3.97/4.00}{Expected May 2028}
\section{AI / ML Experience}
\e{Boston University -- Prof. Brian Kulis}{Boston, MA}{Undergraduate Machine Learning Researcher}{Aug 2026 -- Present}\bi
\item Researching metric learning for human-motion retrieval from temporal 2D/3D pose sequences, adapting contextual-similarity objectives and learning fixed-length embeddings for nearest-neighbor search.
\item Evaluating contextual and contrastive objectives against triplet, supervised-contrastive, and classification retrieval baselines under pose noise, occlusion, temporal jitter, and viewpoint shift using Recall@K and mAP.\ei
\e{Banca Mediolanum}{Milan, Italy}{Software / AI Intern, Customer Knowledge}{Jun 2026 -- Aug 2026}\bi
\item Led development of an internal LLM-agent platform on Databricks from architecture through deployment, integrating governed SQL tools, document RAG, evidence-grounded review, tracing, evaluation, and enterprise chat workflows.
\item Built the agent with MLflow ResponsesAgent and LangGraph ReAct, including versioned Pydantic contracts, request-scoped context, policy-controlled tool execution, citation-backed retrieval, and MLflow tracing.
\item Built prompt-governance and evaluation infrastructure with versioned Delta prompts, side-by-side benchmarks, LLM judging, MLflow GenAI scoring, and Unity Catalog quality controls.\ei
\e{Boston University College of Engineering}{Boston, MA}{Undergraduate Machine Learning Researcher}{Jan 2026 -- May 2026}\bi
\item Built a Python computer-vision pipeline using MMPose and MotionBERT for video stabilization, 2D pose inference, 3D skeleton reconstruction, and stroke-level kinematic analysis.
\item Trained sequence-to-sequence regression models mapping 3D kinematics to rowing force curves and validated predictions against instrumented on-water and erg telemetry.\ei
\section{Engineering Experience}
\e{Societ\`a Cappelletto S.r.l.}{Treviso, Italy}{Software Engineer}{May 2025 -- Oct 2025}\bi\item Rebuilt an Electron inventory application in Tauri, React, and Rust, reducing app size 70\%, memory 80\%, startup time 60\%, and Shopify search latency 43\%.\ei
\e{TickIT GmbH}{Remote}{Frontend \& Backend Engineer}{Nov 2024 -- Feb 2025}\bi\item Shipped Ruby and Next.js/TypeScript ticketing, access-control, analytics, and forecasting features as one of four developers on a 150,000-line platform.\ei
\section{Selected Project}
\textbf{NoteWorthy -- Automated Note Conversion Platform}\\[-2pt]\bi\item Built a Next.js/TypeScript application converting handwritten notes into styled PDFs and LaTeX using Google AI Studio, OAuth, a Dockerized LaTeX compiler, Google Cloud Run, and Vercel.\ei
\section{Leadership}\textbf{Boston University Men's Rowing -- NCAA Division I Student-Athlete} \hfill 2024 -- Present\\Freshman Student-Athlete of the Year (2025); Most Improved Oarsman (2026).
\section{Skills}\textbf{Languages:} Python, TypeScript/JavaScript, SQL, Rust, C/C++\\\textbf{AI/ML:} metric learning, LLM agents, RAG, MLflow, LangGraph, TensorFlow/Keras, MMPose, MotionBERT\\\textbf{Data/Cloud:} Databricks, PostgreSQL, Supabase, Docker, Google Cloud Run, Git
\end{document}
